\documentclass[Orbiter Technical Reference.tex]{subfiles}
\begin{document}

\section{Mesh-to-mesh collision detection and response}
\textbf{EvESpirit}\\
\textbf{May 13, 2026}


\subsection{Introduction}
Orbiter newly implements a custom, yet using already pre-established techniques demonstrated with ~\cite{moller1997} and~\cite{ibarra2003} papers,
mesh-accurate, rigid-body based collision system for vessel-to-vessel,
vessel-to-base, and vessel-to-scatter (foilage or rocks - part of a new Scatterer system) interactions.
Collision remains optionally toggleable for all three categories within Launchpad -> Options -> Physics settings.

Collision detection uses a point-to-triangle intersection algorithm that is
mathematically equivalent to the
M\"oller-Trumbore ray-triangle intersection test~\cite{moller1997},
but solved for a stationary point rather than a directed ray.
Where M\"oller-Trumbore asks ``does this ray hit the triangle?'', the
point-to-triangle variant asks ``is this point inside the triangle's
projection?''.
The collision \emph{response} is handled by a standard Newtonian rigid-body
impulse solver~\cite{ibarra2003}, which computes normal and tangential
impulses and applies them as linear and angular velocity changes.

To maintain framerates that are relatively equivalent
or only mildly affected within the simulation, the system is split into two phases:
a \emph{broad phase} that cheaply discards non-colliding pairs, and a
\emph{narrow phase} that performs the full geometric intersection test.


\subsection{Broad phase - bounding-sphere rejection}
\label{ssec:broadphase}
Given two vessels $A$ and $B$ with bounding-sphere radii
$r_A$ and $r_B$, the broad phase operates on their relative position and
velocity vectors:
\begin{equation*}
\vec{p}_{\mathrm{rel}} = \vec{p}_B - \vec{p}_A, \qquad
\vec{v}_{\mathrm{rel}} = \vec{v}_B - \vec{v}_A
\end{equation*}
The combined bounding radius is $r_{\Sigma} = r_A + r_B$.
A pair is immediately discarded if the vessels are separating
\emph{and} already well apart:
\begin{equation*}
\vec{p}_{\mathrm{rel}} \cdot \vec{v}_{\mathrm{rel}} \ge 0
\;\;\text{and}\;\;
|\vec{p}_{\mathrm{rel}}| > r_{\Sigma} + \epsilon
\end{equation*}
where $\epsilon$ is a small safety margin (5\,m in the initial implementation).

The 5\,m value was chosen to be larger than typical vessel surface 
irregularities ($\sim$1--2\,m for detailed meshes) while remaining small 
enough to avoid premature engagement at long range. This prevents 
false negatives when vessels have complex geometry that extends beyond 
their geometric centers, while maintaining tight broad-phase culling.

For approaching vessels ($\vec{p}_{\mathrm{rel}} \cdot \vec{v}_{\mathrm{rel}} < 0$),
the time of closest approach is estimated as
\begin{equation*}
t_{\mathrm{min}} = -\frac{\vec{p}_{\mathrm{rel}} \cdot \vec{v}_{\mathrm{rel}}}
                         {|\vec{v}_{\mathrm{rel}}|^2}
\end{equation*}
If $t_{\mathrm{min}}$ falls within the current simulation time step $\Delta t$,
the predicted closest-approach distance is checked against $r_{\Sigma}$:
\begin{equation*}
|\vec{p}_{\mathrm{rel}} + \vec{v}_{\mathrm{rel}}\,t_{\mathrm{min}}|
> r_{\Sigma} + \epsilon
\;\;\Rightarrow\;\; \text{skip pair}
\end{equation*}
Only pairs that survive both tests proceed to the narrow phase.


\subsection{Narrow phase - point-to-triangle intersection}
\label{ssec:narrowphase}

\subsubsection{Coordinate transformation}
A cached \emph{point cloud} of hull vertices from vessel~$A$ is transformed
into the local mesh space of vessel~$B$ using the composite rotation matrix
$\mat{R} = \mat{R}_{\mathrm{planet}}^{-1}\,\mat{R}_{\mathrm{vessel}}$.
The inverse transformation is also applied (hull of $B$ tested against
meshes of~$A$) for symmetric detection.

\subsubsection{AABB culling}
Before any per-triangle test, an Axis-Aligned Bounding Box (AABB) overlap
check is performed at two levels:

\begin{enumerate}
\item \textbf{Mesh group AABB.} The AABB of each mesh group in the target
      vessel is compared against the AABB of the incoming hull point cloud.
      Non-overlapping groups are skipped entirely.
\item \textbf{Per-triangle AABB.} Each triangle's bounding box is tested
against individual hull points. A small tolerance of 0.1\,m is added
to both AABBs to account for numerical precision loss in the 
coordinate transformations between vessel-local and mesh-local spaces.
Without this tolerance, edge-case contacts (vertices exactly on AABB
faces etc.) would be incorrectly culled due to floating-point rounding.
\end{enumerate}

\subsubsection{Point-in-triangle test via barycentric coordinates}
\label{sssec:barycentric}
For each triangle with vertices $\vec{p}_0$, $\vec{p}_1$, $\vec{p}_2$
and a candidate hull point $\vec{q}$, define the edge vectors
\begin{equation*}
\vec{e}_1 = \vec{p}_1 - \vec{p}_0, \qquad
\vec{e}_2 = \vec{p}_2 - \vec{p}_0
\end{equation*}
and the triangle face normal
\begin{equation*}
\hat{\vec{n}} = \frac{\vec{e}_1 \times \vec{e}_2}{|\vec{e}_1 \times \vec{e}_2|}
\end{equation*}
The signed distance from the point to the triangle plane is
\begin{equation}
d = (\vec{q} - \vec{p}_0) \cdot \hat{\vec{n}}
\label{eq:signed_dist}
\end{equation}
Points with $d > 0.05\,\text{m}$ (outside) or $d < -2.0\,\text{m}$
(too deep to be a contact event) are rejected. The 0.05\,m outer 
threshold provides a predictive window, in which points approaching within this
distance will likely penetrate in the next frame, allowing the broad 
phase to keep the pair active. The -2.0\,m inner limit prevents 
erroneously large impulses when tunneling has already occurred, and once
penetration exceeds vessel-scale dimensions, position correction
(Section~\ref{ssec:response}) is more appropriate than impulse response.
Additionally, triangles whose normals face away from the relative velocity
vector ($\vec{v}_{\mathrm{rel}} \cdot \hat{\vec{n}} > 0.1$) are skipped.
The 0.1\,m/s threshold (rather than strict $> 0$) prevents oscillation
when vessels are in sliding contact with nearly-tangential velocities;
strictly back-facing surfaces ($\vec{v} \cdot \hat{n} \gg 0.1$) can be
safely ignored, while surfaces within the threshold are kept active to
handle impacts.

The barycentric coordinates $(u, v, w)$ of the projected point are computed
via the Gram matrix of the edge vectors:
\begin{equation}
\begin{aligned}
d_{00} &= \vec{e}_1 \cdot \vec{e}_1, \quad
d_{01} = \vec{e}_1 \cdot \vec{e}_2, \quad
d_{11} = \vec{e}_2 \cdot \vec{e}_2 \\[4pt]
d_{20} &= (\vec{q} - \vec{p}_0) \cdot \vec{e}_1, \quad
d_{21} = (\vec{q} - \vec{p}_0) \cdot \vec{e}_2 \\[4pt]
D &= d_{00}\,d_{11} - d_{01}^{\,2}
\end{aligned}
\label{eq:gram}
\end{equation}
The barycentric weights are then
\begin{equation}
v = \frac{d_{11}\,d_{20} - d_{01}\,d_{21}}{D}, \qquad
w = \frac{d_{00}\,d_{21} - d_{01}\,d_{20}}{D}, \qquad
u = 1 - v - w
\label{eq:bary}
\end{equation}
A hit is registered when $u \ge -0.05$, $v \ge -0.05$, $w \ge -0.05$.
The small negative tolerance (-0.05 $\approx$ 5\% of triangle edge length for
typical meshes) accommodates two sources of error: (i) floating-point
imprecision in the barycentric computation when points lie exactly on
triangle edges, and (ii) slight coordinate misalignment from the
dual-direction testing (hull A → mesh B AND hull B → mesh A).
Without this tolerance, edge and vertex contacts would intermittently fail
detection.
The system tracks the deepest penetration depth $-d$ across all tested
point-triangle pairs and records the corresponding surface normal and
contact location.


\subsection{Collision response - rigid-body impulse solver}
\label{ssec:response}

\subsubsection{Contact velocity}
Once the deepest penetration point and surface normal $\hat{\vec{n}}$
are known, the contact-point velocity for each vessel is computed
including both linear and angular contributions:
\begin{equation}
\vec{v}_A^c = \vec{v}_A + \vec{\omega}_A \times \vec{r}_A, \qquad
\vec{v}_B^c = \vec{v}_B + \vec{\omega}_B \times \vec{r}_B
\label{eq:contact_vel}
\end{equation}
where $\vec{r}_A$ and $\vec{r}_B$ are the moment arms from each vessel's
centre of mass to the contact point.
The relative contact velocity is
\begin{equation*}
\vec{v}_{\mathrm{rel}}^c = \vec{v}_A^c - \vec{v}_B^c
\end{equation*}
and the component along the collision normal is
\begin{equation*}
v_n = \vec{v}_{\mathrm{rel}}^c \cdot \hat{\vec{n}}
\end{equation*}
If $v_n > 0$ the bodies are already separating and no impulse is applied.
If $|v_n| > 500\,\text{m/s}$ a hypervelocity impact is declared and
both vessels are destroyed.

\subsubsection{Normal impulse magnitude}
The impulse magnitude $j$ along the collision normal is computed as
\begin{equation}
j = \frac{-(1+e)\,v_n}
     {\dfrac{1}{m_A} + \dfrac{1}{m_B}
      + \bigl[(\mat{I}_A^{-1}(\vec{r}_A \times \hat{\vec{n}})) \times \vec{r}_A
        + (\mat{I}_B^{-1}(\vec{r}_B \times \hat{\vec{n}})) \times \vec{r}_B\bigr]
        \cdot \hat{\vec{n}}}
\label{eq:impulse}
\end{equation}
where $e$ is the coefficient of restitution (0.2 for vessel-to-vessel,
variable for vessel-to-base), $m_A$, $m_B$ are the vessel masses, and
$\mat{I}_A^{-1}$, $\mat{I}_B^{-1}$ are the inverse inertia tensors
expressed in the global frame.
Because the system evaluates only a single deepest-penetration point
per frame, the inverse inertia values are scaled by a factor of $0.5$
to approximate multi-point contact and reduce rotational jitter.

The denominator in Eq.~\ref{eq:impulse} can be decomposed for clarity.
Define the \emph{angular effect} term:
\begin{equation}
\alpha = \bigl[(\mat{I}_A^{-1}\,\vec{c}_A) \times \vec{r}_A
         + (\mat{I}_B^{-1}\,\vec{c}_B) \times \vec{r}_B\bigr]
         \cdot \hat{\vec{n}}
\label{eq:angular_effect}
\end{equation}
where $\vec{c}_A = \vec{r}_A \times \hat{\vec{n}}$ and
$\vec{c}_B = \vec{r}_B \times \hat{\vec{n}}$.
Then Eq.~\ref{eq:impulse} simplifies to
\begin{equation}
j = \frac{-(1+e)\,v_n}{\dfrac{1}{m_A} + \dfrac{1}{m_B} + \alpha}
\label{eq:impulse_simple}
\end{equation}

\subsubsection{Applying the normal impulse}
The velocity and angular velocity changes are
\begin{equation}
\begin{aligned}
\Delta\vec{v}_A &= +\frac{j\,\hat{\vec{n}}}{m_A}, &\qquad
\Delta\vec{v}_B &= -\frac{j\,\hat{\vec{n}}}{m_B} \\[4pt]
\Delta\vec{\omega}_A &= +\mat{I}_A^{-1}\,\vec{c}_A\;j, &\qquad
\Delta\vec{\omega}_B &= -\mat{I}_B^{-1}\,\vec{c}_B\;j
\end{aligned}
\label{eq:apply_impulse}
\end{equation}

\subsubsection{Tangential friction}
The tangential component of the relative contact velocity is
\begin{equation*}
\vec{v}_t = \vec{v}_{\mathrm{rel}}^c - v_n\,\hat{\vec{n}}
\end{equation*}
A friction impulse is computed using the same impulse denominator structure
as the normal case but along the tangent direction
$\hat{\vec{t}} = \vec{v}_t / |\vec{v}_t|$:
\begin{equation}
j_t = \min\!\bigl(|\vec{v}_t|\,m_{\mathrm{eff}}^t,\;\mu\,j\bigr)
\label{eq:friction}
\end{equation}
where $\mu$ is the Coulomb friction coefficient (0.6 for vessel-to-vessel)
and $m_{\mathrm{eff}}^t$ is the effective mass in the tangential direction,
computed analogously to the normal case.
The friction impulse is capped by $\mu\,j$ (Coulomb's law) to prevent
the tangential correction from reversing the sliding direction.

For vessel-to-base collisions, wheel brakes are modelled by scaling $\mu$
from a base value of 0.8 up to 2.8 at full brake application:
\begin{equation*}
\mu = 0.8 + 2.0\,b
\end{equation*}
where $b \in [0,1]$ is the average of left and right wheel brake levels.

\subsubsection{Positional correction}
After all impulses are applied, a direct positional correction
pushes the interpenetrating bodies apart:
\begin{equation}
\begin{aligned}
\Delta\vec{p}_A &= +\hat{\vec{n}}\,d_{\mathrm{pen}} \cdot 1.01 \cdot
                    \frac{m_B}{m_A + m_B} \\[4pt]
\Delta\vec{p}_B &= -\hat{\vec{n}}\,d_{\mathrm{pen}} \cdot 1.01 \cdot
                    \frac{m_A}{m_A + m_B}
\end{aligned}
\label{eq:pos_correction}
\end{equation}
where $d_{\mathrm{pen}}$ is the measured penetration depth.
The 1\% over-correction ($\times 1.01$) prevents persistent
frame-to-frame interpenetration in two scenarios: (i) when vessels
remain in contact across multiple frames (resting contact), accumulated
floating-point drift can allow slow interpenetration growth - the slight
push ensures separation; (ii) when the collision normal is nearly 
tangent to curved surfaces, exact separation may leave points marginally
inside the next adjacent triangle. The 1\% value was empirically tuned
to eliminate jitter (smaller values allow persistent contact) without
causing visible "bouncing" separation artifacts (larger values would).


\subsection{Vessel-to-base collisions}
\label{ssec:vessel_base}
Vessel-to-base collisions follow the same narrow-phase point-to-triangle
algorithm described in Section~\ref{ssec:narrowphase}.
The base is treated as an immovable object, so the impulse solver
simplifies: $m_B \to \infty$ and $\mat{I}_B^{-1} = \mat{0}$.
The normal impulse reduces to
\begin{equation*}
j_n = -(1 + e)\,v_n
\end{equation*}
with $e = 0.1$ for impacts exceeding 0.5\,m/s and $e = 0$ for softer
contacts to prevent resting-state jitter.

Angular damping is applied when a vessel is nearly at rest on a surface
($|v_n| < 0.5\,\text{m/s}$ and $d_{\mathrm{pen}} > 0.005\,\text{m}$):
the angular velocity relative to the planet's surface rotation is reduced
by 20\% per frame. This compensates for the single-contact-point
approximation and prevents tipping oscillations by 20\% per frame.
The 0.5\,m/s velocity threshold distinguishes  "resting" from "impacting" states,
as below this speed, contacts are  treated as static friction scenarios requiring stabilization; above it,
normal restitution physics apply. The 0.005\,m penetration requirement
ensures damping only activates when meaningful surface contact exists,
preventing false activation during grazing trajectories. The 20\% 
damping rate (rather than aggressive 50%+ or conservative 5%) balances
two competing needs: fast convergence to rest (achieved in $\sim$10 frames)
versus smooth, artifact-free motion during slow rolling.


\subsection{Base configuration - the COLLISION marker}
\label{ssec:cfg_collision}
Base configuration files (\texttt{.cfg}) can optionally flag individual
mesh entries for collision detection using the \texttt{COLLISION} keyword.
If a mesh block does not contain this marker, no collision check is
performed against it (preserving legacy behaviour).

\begin{lstlisting}[language=OSFS]
HANGAR
  POS 0 0 0
  ROT 90
  SCALE 12 8 244
  TEX1 concrete1 0.4 0.4
  TEX3 concrete1 0.5 8
  COLLISION
END
\end{lstlisting}

The \texttt{COLLISION} keyword can appear anywhere within the mesh block
between its header and the closing \texttt{END} statement.
Multiple mesh entries in the same base configuration may independently
use or omit the marker.

\end{document}
